\documentclass[11pt]{article}
\usepackage[T1]{fontenc}
\usepackage{lmodern}
\usepackage[utf8]{inputenc}
\usepackage{amsmath,amssymb,amsthm,mathtools}
\usepackage{microtype}
\usepackage[margin=1in]{geometry}
\usepackage{enumitem}
\usepackage{hyperref}
\hypersetup{colorlinks=true,linkcolor=blue,citecolor=blue,urlcolor=blue}

\newtheorem{theorem}{Theorem}
\newtheorem{proposition}[theorem]{Proposition}
\newtheorem{corollary}[theorem]{Corollary}
\newtheorem{hypothesis}[theorem]{Hypothesis}
\theoremstyle{definition}
\newtheorem{definition}[theorem]{Definition}
\newtheorem{convention}[theorem]{Convention}
\newtheorem{remark}[theorem]{Remark}
\newtheorem{example}[theorem]{Example}

\newcommand{\dR}{\mathrm{dR}}
\newcommand{\Omop}{\Omega_{\mathrm{op}}}
\newcommand{\Omres}{\Omega^{\mathrm{res}}}
\newcommand{\Sigspin}{\Sigma_{\mathrm{spin}}}
\DeclareMathOperator{\opskew}{skew}

\title{Two Channels, Three Readouts:\\
A Typing Refinement of Structural Spin}
\author{Maksim Barziankou}
\date{July 2026}

\begin{document}
\maketitle

\begin{abstract}
Structural spin is the non-gradient layer \(S\) in a flow
\(\dot x=-\nabla V+S\).  Its projection-relative signature has two formally
independent components: a local differential channel \(d(S^\flat)\), whose
half is operational spin, and a global cohomological channel
\([\eta_S]\in H^1_{\dR}(M)\), whose periods provide winding witnesses.
This note separates that two-channel structure from three scalar readout
families---period, path norm/amplitude, and vorticity flux---and from a
fourth object, the dynamic residual used in drift or cost ledgers.

The distinction is load-bearing.  A non-zero period forces a positive
path norm when both are read from the same declared one-form; this is
the triangle inequality, not a lower bound on vorticity.  Conversely,
neither a path norm nor a vorticity flux determines the cohomology class.
Explicit local-only and topological-only examples show that the two
channels are logically independent over the class of all declared signature
data, while an equal-norm collision on a flat torus shows that an
amplitude-only monitor is incomplete for winding.  On the compact
admissible branch, the topological-only sector remains premise-gated by
the harmonic obstruction space \(W_V\).  A registered finite spanning
family of harmonic-evaluation covectors reduces this gate exactly to the
zero test \(Ea=0\), so algebraic independence is not confused with
unrestricted realizability.

The resulting typing principle requires every bridge from structural
spin to a measured rotation index, dynamic residual, or maintenance cost
to declare its probe and observation map.  Applied to cortical rotating
waves, it yields a narrow empirical hypothesis: winding can collapse
while \(2\text{--}8\) Hz amplitude remains above a calibrated detection
floor and within a predeclared equivalence margin.  The dual local-only arm uses an antisymmetric-Jacobian score defined
without explicitly consuming the registered amplitude functional, a positive
instrument/effect floor, and an exchangeability-calibrated rank threshold.
The stronger \emph{survives-but-dead} verdict additionally requires both an
independent viability/continuation control and a predeclared
identity-witness bridge.
\end{abstract}

\bigskip
\noindent\footnotesize
\textbf{License \& attribution.}\ \copyright~2026 Maksim Barziankou.
Licensed under
\href{https://creativecommons.org/licenses/by-nc-nd/4.0/}{CC\,BY-NC-ND\,4.0}.
\par\smallskip
\noindent\textbf{DOI.}\ \href{https://doi.org/10.17605/OSF.IO/69FE2}{10.17605/OSF.IO/69FE2}.
\par\smallskip
\noindent No patent license is granted.  Part of the
\emph{Navigational Cybernetics~2.5} research programme.  This is a
standalone companion to \emph{Structural Spin as a Forgetful Separation
over Dissipative Dynamics, with a Cohomological Obstruction to Identity
Transfer}, DOI:10.17605/OSF.IO/94GWQ.
\normalsize
\bigskip

\section{Purpose and relation to the preceding work}

The preceding structural-spin paper established a faithful-not-full
forgetful separation over a gradient skeleton, a cohomological witness
for topological identity transfer, the necessity/blindness limits of that
witness, and a compact-branch obstruction to eliminating operational
spin while preserving an admissible non-zero class
\cite{structural}.  It also stated the central decomposition: one
non-gradient layer has a local vorticity channel and a global
cohomological channel.

The present note does not replace those results.  It resolves a narrower
typing ambiguity created when the phrases \emph{local channel},
\emph{spin magnitude}, \emph{vorticity flux}, and \emph{residual spin}
are allowed to stand for one another.  They do not have the same type.
The correction can be stated in one line:
\[
\boxed{\begin{gathered}
\text{one carrier}\ \longrightarrow\ \text{two structural channels}\\[-1pt]
\text{three readout families}\ \longrightarrow\ \text{an explicitly bridged dynamic ledger}
\end{gathered}}
\]

This separation sharpens the spin-channel bridge \cite{bridge}, the
physical operational-spin construction \cite{operational}, and the
consolidated NC2.5 programme statement \cite{nc25}.  It also isolates
exactly what a rotating-wave
experiment can test without importing the identity semantics of the
abstract theory.

\section{Carrier and structural signature}

Let \((M,g)\) be an oriented Riemannian manifold, let \(\nabla\) denote its
Levi-Civita connection, and let
\[
F=-\nabla V+S,\qquad \omega_S:=S^\flat.
\]
The metric is fixed; the musical isomorphism and all norm statements are
relative to \(g\).  A splitting \(F=-\nabla V+S\) supplies a representative,
but the signature below is invariant under its gauge once the selector is
fixed: the substitution \((V,S)\mapsto(V+u,S+\nabla u)\) leaves \(F\) unchanged
while replacing \(\omega_S\) by \(\omega_S+du\).  Throughout,
the exterior derivative of a \(1\)-form uses the convention
\[
d\omega(X,Y)=X\omega(Y)-Y\omega(X)-\omega([X,Y]),
\]
so that \(d(-y\,dx+x\,dy)=2\,dx\wedge dy\); the factor \(\tfrac12\) in
Definition~\ref{def:signature} is stated relative to this convention and
changes under the alternating normalization.  On a compact oriented
boundaryless manifold, Hodge theory \cite{warner} gives
\[
\omega_S=df+\delta\beta+h.
\]
On the alternative closed-component branch, assume a fixed declared linear
projection
\[
P_{\mathrm{cl}}:\Omega^1(M)\longrightarrow Z^1(M),\qquad
Z^1(M):=\ker\!\left(d:\Omega^1(M)\to\Omega^2(M)\right),
\]
and set \(\eta_S:=P_{\mathrm{cl}}(\omega_S)\).  Require
\(P_{\mathrm{cl}}\circ P_{\mathrm{cl}}=P_{\mathrm{cl}}\), and require the map
to preserve cohomology on closed inputs: for every closed \(\alpha\in
\Omega^1(M)\), the form \(P_{\mathrm{cl}}(\alpha)-\alpha\) is exact.  (The
harmonic projection satisfies both; note that it is \emph{not} the identity
on \(Z^1(M)\), so idempotence may not be strengthened to
\(P_{\mathrm{cl}}|_{Z^1}=\mathrm{id}\).)  These conditions constrain \(P_{\mathrm{cl}}\) only on closed inputs.  They
nevertheless make the splitting gauge harmless: linearity and closed-input
cohomology preservation give
\([P_{\mathrm{cl}}(du)]=[du]=0\), while
\(d(\omega_S+du)=d\omega_S\).  Thus the signature below is an invariant of
the declared pair \((F,P_{\mathrm{cl}})\), even though its representative
\(\eta_S\) may change.  The selector remains part of the type: when
\(\omega_S\) is not closed, different admissible selectors may assign
different classes to the same \(F\), and those classes may not be compared.
In the compact case one may take
\(\eta_S=h\), the harmonic projection, which satisfies the requirement:
for closed \(\alpha=df+\delta\beta+h\) one has \(d\delta\beta=0\), since
\(d\alpha=0\) and \(ddf=dh=0\), whence
\(\lVert\delta\beta\rVert^2=\langle\beta,d\delta\beta\rangle=0\),
\(\delta\beta=0\) and \(\alpha-h=df\) is exact.

De Rham cohomology is used throughout in the sense of \cite{bott-tu}.

\begin{definition}[Structural-spin signature]\label{def:signature}
On the declared-selector branch, the \emph{structural-spin signature} is
\[
\Sigspin(F;P_{\mathrm{cl}})
  :=\bigl([\eta_S],d\omega_S\bigr)
  \in H^1_{\dR}(M)\times\Omega^2(M).
\]
On the compact Hodge branch the harmonic projection is canonically fixed by
\(g\); there we suppress it and write \(\Sigspin(F)\).  Selector arguments
may be suppressed nowhere else.  The \emph{topological channel} is
\([\eta_S]\).  The \emph{local differential channel} is \(d\omega_S\),
with operational spin
\[
\Omop:=\frac12d\omega_S.
\]
\end{definition}

Tensorially, let \(K_{\mathrm{op}}\in\operatorname{End}(TM)\) be the
skew-adjoint endomorphism defined by
\[
g(K_{\mathrm{op}}X,Y)=\Omop(X,Y)=\frac12d\omega_S(X,Y).
\]
If \(J_F(X):=\nabla_XF\) and \(J_F^{*g}\) is its metric adjoint, then
\[
K_{\mathrm{op}}=\frac12(J_F-J_F^{*g}).
\]
Indeed, metric compatibility gives \((\nabla_X\omega_S)(Y)=g(\nabla_XS,Y)\)
and torsion-freeness gives
\(d\omega_S(X,Y)=(\nabla_X\omega_S)(Y)-(\nabla_Y\omega_S)(X)\), so that
\(d\omega_S(X,Y)=g\bigl((J_S-J_S^{*g})X,Y\bigr)\) for \(J_S(X):=\nabla_XS\).
Torsion-freeness also makes \(\nabla^2V\) symmetric, so the gradient layer is
self-adjoint and contributes nothing to the skew part:
\(J_F-J_F^{*g}=J_S-J_S^{*g}\).  The local differential channel is therefore
an invariant of \(F\) alone.  The topological channel is also splitting-gauge
invariant for fixed \(P_{\mathrm{cl}}\), but remains selector-relative; hence
the full minimal type is \((F,P_{\mathrm{cl}})\), or \(F\) relative to the
fixed metric on the compact Hodge branch.
Only in an orthonormal frame does this become the matrix formula
\(K_{\mathrm{op}}=\tfrac12(J_F-J_F^\top)\).
The physical construction in \cite{operational} applies the same
antisymmetrization type, \(\opskew(A):=\tfrac12(A-A^{*g})\), to a
directed-change field \(J_t=\partial_tu_t\):
\(\Omega_t=\opskew(\nabla J_t)\).  The fields \(F\) and \(J_t\) are not
thereby identified; what is shared is the operation and its tensor type.

\begin{remark}[Premise-minimal scope]
Architectural spin can be asserted more weakly as an obstruction to a
global scalar-gradient realization.  Definition~\ref{def:signature}
requires the additional representative/regularity branch.  Without it,
neither component of the displayed pair may be silently manufactured.
\end{remark}

\section{Formal independence of the two channels}

\begin{proposition}[Channel independence]\label{prop:independence}
The two vanishing conditions are logically independent over the class of
all declared signature data---quadruples \((M,g,F,P_{\mathrm{cl}})\) on the
selector branch and triples \((M,g,F)\) on the compact Hodge branch: neither
\[
d\omega_S\neq0\ \Longrightarrow\ [\eta_S]\neq0
\qquad\text{nor}\qquad
[\eta_S]\neq0\ \Longrightarrow\ d\omega_S\neq0
\]
holds universally over that class.
\end{proposition}

\begin{proof}
Work on the single flat torus
\(M=\mathbb R^2/(2\pi\mathbb Z)^2\) with its standard metric and canonical
Hodge projection.  Put \(V=0\), so \(F=S\), for both witnesses.  For the
local-only direction take
\(S_{\mathrm L}=\sqrt2\cos y\,\partial_x\).  Then
\[
\omega_{\mathrm L}=\sqrt2\cos y\,dx,\qquad
d\omega_{\mathrm L}=\sqrt2\sin y\,dx\wedge dy\neq0,
\]
while the harmonic projection of \(\omega_{\mathrm L}\) is zero because its
coefficient has zero torus mean.  Thus \([\eta_{\mathrm L}]=0\).

For the topological-only direction on the same \((M,g)\), take
\(S_{\mathrm T}=\partial_x\).  Then
\[
\omega_{\mathrm T}=dx,\qquad d\omega_{\mathrm T}=0,\qquad
\oint_{\{y=0\}}dx=2\pi.
\]
Since \(dx\) is harmonic, \([\eta_{\mathrm T}]=[dx]\neq0\), while the local
differential channel vanishes.  Both vector fields are non-gradient and
divergence-free; divergence-freeness is an additional property of these
witnesses, not part of the definition of non-gradient.
\end{proof}

\begin{corollary}[No channel substitution]\label{cor:no-sub}
No inference from the non-vanishing of one channel to the non-vanishing of
the other is licensed by the signature alone: a local vorticity detector is
not a complete topological-witness detector, and a winding detector is not a
complete local-vorticity detector.  In particular, the phrase \emph{two
channels} refers to \(([\eta_S],d\omega_S)\), not to winding and amplitude.
\end{corollary}

\begin{remark}[Scope of the independence claim]
The proof realizes both one-channel directions on one fixed compact
\((M,g)\), on the canonical Hodge branch.  The same torus also realizes the
jointly non-zero pattern with
\(\omega=(1+\sqrt2\cos y)\,dx\).  A zero--zero form such as
\(d(\cos x)\) is exact and therefore gradient; it lies outside the declared
non-gradient carrier class.  Hence all three admissible non-zero sign patterns
occur on this one compact manifold, without a claim that a fourth
non-gradient zero--zero pattern exists.
\end{remark}

\section{Three readout families}

Structural channels are not yet scalar measurements.  A measurement
requires a probe.

\begin{definition}[Typed readouts]\label{def:readouts}
Let \(\gamma:[0,T]\to M\) be a piecewise \(C^1\) closed curve.  Let
\(\mathcal P\) be an oriented compact piecewise \(C^1\) \(2\)-chain
(or embedded probe surface, possibly with boundary) on which the
pullback of \(d\omega_S\) is integrable.  For a declared one-form
\(\alpha\), define
\begin{align}
W_\gamma(\alpha)
  &:=\oint_\gamma\alpha,
  &&\text{period/winding readout},\label{eq:period}\\
A_\gamma(\alpha)
  &:=\int_0^T
       \bigl|\alpha_{\gamma(t)}(\dot\gamma(t))\bigr|\,dt,
  &&\text{path-amplitude/total-variation readout},\label{eq:norm}\\
Q_{\mathcal P}(\omega_S)
  &:=\int_{\mathcal P}d\omega_S,
  &&\text{oriented vorticity-flux readout}.\label{eq:flux}
\end{align}
Here \(W_\gamma\) and \(A_\gamma\) are functionals of \(\alpha\), and
\(Q_{\mathcal P}\) is a functional of \(\omega_S\); their displayed
evaluated values are real scalars.  When \(\alpha=\eta_S\) is closed,
\(W_\gamma(\eta_S)\) depends only on the homology class of \(\gamma\).
The readout \(A_\gamma\) is invariant under orientation-preserving
reparametrization, by the change of variables
\(|\alpha(\dot\gamma\circ\varphi)\varphi'|\,ds=|\alpha(\dot\gamma)|\,dt\),
so it depends only on the traversed path and not on the chosen
\([0,T]\); it is a \emph{seminorm} in \(\alpha\), since
\(A_\gamma(\alpha)=0\) whenever \(\alpha\) annihilates \(\dot\gamma\)
along \(\gamma\), without \(\alpha=0\).
All probes and geometries remain declared.  If \(\mathcal P\) is a
closed \(2\)-cycle and \(\omega_S\) is globally defined on it, then
\(Q_{\mathcal P}(\omega_S)=0\) by Stokes.  Throughout this note
\(\omega_S\) is globally defined, so a non-zero flux readout is used only
for a chain with \(\partial\mathcal P\neq0\), where
\[
Q_{\mathcal P}(\omega_S)=W_{\partial\mathcal P}(\omega_S).
\]
Here \(W\) is extended additively over the components of the boundary
cycle.  Patchwise connections, distributional singularities and removed-core
fluxes require a separate carrier definition and are outside the present
scope.  Stokes is the one determinate relation holding between two of these
readouts without further model data.
\end{definition}

\begin{proposition}[The norm inequality is not a vorticity
inequality]\label{prop:norm}
For every declared pair \((\gamma,\alpha)\),
\[
A_\gamma(\alpha)\ge |W_\gamma(\alpha)|.
\]
If \(\alpha=\eta_S\), a non-zero topological period therefore forces a
positive path norm of that representative.  It does not force
\(d\omega_S\neq0\), and it supplies no lower bound on
\(|Q_{\mathcal P}(\omega_S)|\).
\end{proposition}

\begin{proof}
The displayed inequality is
\[
\int_0^T|\alpha(\dot\gamma)|\,dt
 \ge
\left|\int_0^T\alpha(\dot\gamma)\,dt\right|,
\]
the ordinary triangle inequality.  The topological-only member of
Proposition~\ref{prop:independence} has non-zero period and positive path
norm but \(d\omega_S=0\), hence every vorticity-flux readout is zero.
\end{proof}

\begin{remark}[Instrument compatibility]
If an experiment measures amplitude from a signal other than the
one-form used for the period, Proposition~\ref{prop:norm} does not apply
until a reconstruction map between those observables is declared.  The
same name \emph{rotation} is not such a map.
\end{remark}

\begin{proposition}[An amplitude-only collision]\label{prop:collision}
A monitor reading only the path amplitude of the carrier \(\omega_S\)
cannot recover the topological channel, even on a fixed probe manifold and
loop, and even on the compact Hodge branch: two divergence-free,
non-gradient carriers with equal amplitude readouts can have topological
classes \(0\) and \([dx]\neq0\).
\end{proposition}

\begin{proof}
Take the flat torus
\(M=(\mathbb R/2\pi\mathbb Z)^2\), with coordinates \((x,y)\), and the
oriented loop
\[
\gamma(t)=(t,\pi/4),\qquad 0\le t\le2\pi.
\]
Let
\[
\omega_1=dx,\qquad
\omega_0=\sqrt2\cos y\,dx.
\]
Along \(\gamma\), both forms evaluate to \(1\) on \(\dot\gamma\), hence
\[
A_\gamma(\omega_1)=2\pi=A_\gamma(\omega_0).
\]
The harmonic \(1\)-forms on the flat torus are spanned by \(dx,dy\).
The harmonic projection of \(\omega_1\) is \(dx\), whereas that of
\(\omega_0\) is zero: its \(dx\)-coefficient has zero torus average and its
\(dy\)-coefficient vanishes.  Thus
\[
\eta_1=dx,\qquad \eta_0=0,
\qquad
W_\gamma(\eta_1)=2\pi,\qquad W_\gamma(\eta_0)=0.
\]
Both associated fields are divergence-free.  Moreover
\[
d\omega_0=\sqrt2\sin y\,dx\wedge dy\neq0,
\]
and on \(\gamma\) its coefficient equals \(1\); consequently
\(\omega_0\) is genuinely non-gradient and carries a non-zero local
differential channel despite its zero topological channel.  The collision
therefore remains entirely inside the declared non-gradient carrier class:
an amplitude-only monitor assigns the same value to opposite topological
witness states.
\end{proof}

\begin{remark}[Why the readouts must be typed to their sources]
The collision is between \(A_\gamma(\omega_i)\), read from each carrier,
and \(W_\gamma(\eta_i)\), read from its harmonic component.  It is not a
counterexample to Proposition~\ref{prop:norm}, which compares two readouts
of the \emph{same} declared one-form: for each \(\eta_i\) separately,
\(A_\gamma(\eta_i)\ge|W_\gamma(\eta_i)|\).  The two statements are
compatible precisely because their sources differ, which is the typing
point of Section~\ref{sec:bridge}.
\end{remark}

\section{The dynamic ledger and the typed bridge}\label{sec:bridge}

The residual variable used in a drift law is a fourth object.  Write it
as \(\Omres_t\in\mathcal R_t\), where \(\mathcal R_t\) is the declared
residual space.  A ledger statement has, schematically, the form
\[
D_T\ge c_1\int_0^T\|\Omres_t\|_{\mathcal R_t}\,dt,
\]
where \(D_T\) is the declared accumulated drift or maintenance budget over
\([0,T]\) and \(c_1>0\) is a declared calibration constant with units
\[
[c_1]=\frac{[D_T]}{[\lVert\Omres\rVert]\, [\mathrm{time}]};
\]
neither its value nor its units are fixed by the structural signature.
Nothing in Definition~\ref{def:signature} makes
\(\Omres\) equal to \(d\omega_S\), its flux, or a path norm.

\begin{convention}[Typed bridge principle]\label{thm:typed-bridge}
Let a structural-spin model use any of the following in a measurement,
drift, or cost law:
\[
[\eta_S],\quad d\omega_S,\quad
W_\gamma(\alpha),\quad A_\gamma(\alpha),\quad
Q_{\mathcal P}(\omega_S),\quad \Omres.
\]
Then:
\begin{enumerate}[label=\textup{(\roman*)},leftmargin=2.2em]
  \item the first two are the structural channels;
  \item \(W_\gamma(\alpha),A_\gamma(\alpha)\), and
        \(Q_{\mathcal P}(\omega_S)\) are evaluated probe-dependent
        scalar readouts, not further spin objects;
  \item the period--norm inequality of Proposition~\ref{prop:norm} is
        licensed when both readouts use the same declared one-form and
        locus; that proposition licenses no inequality between readouts of
        \emph{different} declared one-forms, and any such relation requires
        its own derivation;
  \item any relation between \(d\omega_S\) or
        \(Q_{\mathcal P}(\omega_S)\) and \(\Omres\) requires an
        explicit observation map, regulated
        estimator, or sealed dependence;
  \item any scalar cost law must name its readout mode and calibration.
\end{enumerate}
The readouts are related to one another only by Stokes,
\(Q_{\mathcal P}(\omega_S)=W_{\partial\mathcal P}(\omega_S)\), and by
Proposition~\ref{prop:norm}; no equality between a structural channel and
the dynamic residual follows from the structural signature alone.
\end{convention}

\begin{remark}[Justification]
Items (i) and (ii) record type membership; the codomains already separate
the objects:
\[
[\eta_S]\in H^1_{\dR}(M),\quad
d\omega_S\in\Omega^2(M),\quad
W_\gamma(\alpha),A_\gamma(\alpha),Q_{\mathcal P}(\omega_S)
\in\mathbb R,\quad \Omres_t\in\mathcal R_t.
\]
The scalar readouts arise only after evaluation on a curve, normed
locus, or surface.  Item (iii) is a restatement of
Proposition~\ref{prop:norm}; items (iv) and (v) are declared modelling
requirements rather than propositions.  Proposition~\ref{prop:norm} gives
the sole universal period--norm relation used here, and
Proposition~\ref{prop:independence} rules out either structural
non-vanishing implication.  A map such as
\[
\mathcal O_t:\Omega^2(M)\longrightarrow\mathcal R_t
\quad\text{or}\quad
\mathcal B_t:\mathbb R\longrightarrow\mathcal R_t,\qquad
\Omres_t=\mathcal B_t\!\left(Q_{\mathcal P}(\omega_S)\right)
\]
is therefore additional model data.  The same is true of a cost
functional
\[
\mathcal C
 =c_W|W_\gamma(\eta_S)|+c_AA_\gamma(\alpha)
  +c_Q|Q_{\mathcal P}(\omega_S)|,
\]
whose coefficients, units, probes, and active terms must be declared.
\end{remark}

\begin{corollary}[Correct reading of the boundary norm law]
Let \(\sigma_H\) and \(\Lambda_{\mathrm{loc}}\) be the boundary period and
the compatible discrete path norm of \cite{structural}.  The implication
\[
\sigma_H\neq0\quad\Longrightarrow\quad
\Lambda_{\mathrm{loc}}>0
\]
is valid when \(\sigma_H\) is a period and
\(\Lambda_{\mathrm{loc}}\) is the compatible path norm.  It is not an
identity between topological spin and operational spin, and it does not
identify \(\Lambda_{\mathrm{loc}}\) with \(\Omres\).
\end{corollary}

\section{Compact admissibility and the topological-only sector}
\label{sec:compact}

Formal independence does not imply that every pair of channel values is
realizable under every admissibility package.  The compact branch of
\cite{structural} makes this precise.

Let \(M\) be compact and boundaryless and \(V\in C^\infty(M)\).  Write \(\mathcal H^1(M)\) for
the harmonic \(1\)-forms and
\(\Phi:\mathcal H^1(M)\xrightarrow{\sim}H^1_{\dR}(M)\) for the Hodge
isomorphism.  Define
\[
L:\mathcal H^1(M)\to C^\infty(M),\qquad
L(h)=h(\nabla V),
\]
and
\[
W_V:=\Phi(\ker L)\subseteq H^1_{\dR}(M).
\]

\begin{proposition}[Premise gate for topological-only spin]
\label{prop:WV}
Say that \(S\) \emph{carries} the class \(c\) if \([\eta_S]=c\).  A class
\(c\neq0\) is carried by a divergence-free, operational-spin-free field
satisfying \(\langle S,\nabla V\rangle=0\) pointwise if and only if
\(c\in W_V\).  (The equivalence also holds for \(c=0\); the hypothesis
\(c\neq0\) serves only to exclude the trivial witness \(S=0\).)
\end{proposition}

\begin{proof}
Divergence-free and operational-spin-free mean
\(\delta\omega_S=0\) and \(d\omega_S=0\), so
\(\Delta\omega_S=(d\delta+\delta d)\omega_S=0\) and \(\omega_S\) is
harmonic; this step uses neither compactness nor boundarylessness.  Both
are used next, for uniqueness: if a harmonic \(h\) is exact, \(h=d\alpha\),
then \(\lVert h\rVert^2_{L^2}=\langle d\alpha,h\rangle
=\langle\alpha,\delta h\rangle=0\) by integration by parts on a compact
boundaryless \(M\), so \(\mathcal H^1(M)\cap\operatorname{im}d=0\) and each
class has exactly one harmonic representative.  Hence \(\omega_S=h_c\), and
since \(\omega_S\) is already harmonic, \(\eta_S=h=\omega_S\) on this
branch, so \(S\) carries \([\omega_S]\).  Pointwise admissibility is
precisely \(h_c(\nabla V)=0\), equivalently \(h_c\in\ker L\).

Conversely, let \(c\in W_V\), so \(c=\Phi(h)\) with \(h\in\ker L\);
injectivity of \(\Phi\) forces \(h=h_c\).  Put \(S:=h_c^\sharp\).  Then
\(\omega_S=h_c\) is harmonic, so \(d\omega_S=0\) (operational-spin-free)
and \(\delta\omega_S=0\) (divergence-free); moreover
\(\langle S,\nabla V\rangle=h_c(\nabla V)=L(h_c)=0\) pointwise, and
\([\eta_S]=[h_c]=c\).  For \(c\neq0\) we get \(h_c\neq0\), so
\(S\not\equiv0\).
\end{proof}

\begin{corollary}[Finite exact \(W_V\) arm]\label{cor:WV-arm}
For the fixed Riemannian metric and Hodge identification above, let
\(h_1,\ldots,h_b\) be a registered harmonic basis and define
\[
F_V(p)\in\mathcal H^1(M)^*,\qquad F_V(p)(h):=h_p(\nabla V_p).
\]
Let \(p_1,\ldots,p_m\in M\) satisfy the spanning-certificate premise
\[
\operatorname{span}\{F_V(p_j):1\le j\le m\}
=\operatorname{span}\{F_V(p):p\in M\}.
\]
Set \(E_{ji}:=h_i(\nabla V)(p_j)\).  For
\[
c=\Phi\!\left(\sum_{i=1}^b a_i h_i\right),
\qquad a=(a_1,\ldots,a_b)^\top\ne0,
\]
one has
\[
c\in W_V\quad\Longleftrightarrow\quad Ea=0.
\]
Consequently
\[
\rho_V^{(\mathbf h,P)}(c):=\lVert Ea\rVert_\infty,
\qquad \mathbf h=(h_1,\ldots,h_b),\quad P=(p_1,\ldots,p_m),
\]
written \(\rho_V(c)\) whenever the registration \((\mathbf h,P)\) is fixed,
is an executable observable with exact threshold zero when the class
coordinates and matrix entries are exact symbolic or rational data.  A
basis change \(\mathbf h'=\mathbf h T\) gives \(E'=ET\) and
\(a'=T^{-1}a\), hence \(E'a'=Ea\): for a fixed ordered sample tuple the
whole residual vector, and therefore \(\rho_V\), is basis-independent.
Its zero set is also invariant under non-zero row rescaling and replacement
of \(P\) by another certified spanning sample family.  The positive
magnitude does change under row rescaling and sample choice (enlarging
\(P\) can only increase it) and is positively homogeneous of degree one in
\(c\), so it induces no meaningful comparison between classes.  Numerical
harmonic solvers require a separately derived error
bound or interval certificate and are outside this exact arm.  Under the
stated certificates, a declared membership with \(\rho_V(c)>0\), or a
declared non-membership with \(\rho_V(c)=0\), falsifies the registered
claim.  The spanning-certificate premise is load-bearing in exactly one
direction: \(Ea\neq0\) exhibits a point at which \(h_a(\nabla V)\neq0\) and
so excludes membership without any span premise, whereas \(Ea=0\) bounds
\(h_a(\nabla V)\) only on the sampled points.  Accordingly, without that
premise a non-zero sampled residual still excludes membership, but a zero
sampled residual is only non-rejection on the sampled points.
\end{corollary}

\begin{proof}
For \(h_a=\sum_i a_i h_i\), the vector \(Ea\) is
\((F_V(p_1)(h_a),\ldots,F_V(p_m)(h_a))^\top\).  Its vanishing annihilates
the registered span, which equals the span of every \(F_V(p)\); hence
\(h_a(\nabla V)=0\) pointwise.  The reverse implication is immediate.
Proposition~\ref{prop:WV} completes the equivalence.
\end{proof}

\begin{example}[Registered flat-torus instance]\label{ex:WV-torus}
Take
\(M=(\mathbb R/2\pi\mathbb Z)^2\) with
\(g=dx^2+dy^2\), coordinates \((x,y)\), \(V=\cos x\), and harmonic
basis \((dx,dy)\).  Then
\[
L(a\,dx+b\,dy)=-a\sin x.
\]
In the dual basis \(F_V(x,y)=(-\sin x,\,0)\) for every \((x,y)\in M\), so
the whole evaluation family lies on a single line and one point discharges
the spanning certificate: \(p_1=(\pi/2,0)\) gives the spanning row
\(E=(-1,0)\).  (This is a feature of \(\partial_yV\equiv0\), not a general
sufficiency of one sample.)  Thus \(W_V=\operatorname{span}\{[dy]\}\),
consistently with \(\dim W_V=b_1-1=1\), and every
non-zero class \(c_{a,b}=[a\,dx+b\,dy]\) passes the exact gate if and
only if \(a=0\).  This is the registered exact-arithmetic instance used
by the companion code; the code does not authenticate the harmonic-basis
or spanning certificates themselves.
\end{example}

Thus the topological-only sector is real but constrained.  With
\(F_V(p)\) as in Corollary~\ref{cor:WV-arm}, the preceding work proves
\[
\dim W_V
=b_1(M)-\dim_{\mathbb R}
  \operatorname{span}\{F_V(p):p\in M\},
\]
since \(\ker L\) is the annihilator in \(\mathcal H^1(M)\) of that span,
which is finite-dimensional as a subspace of the \(b\)-dimensional dual
\(\mathcal H^1(M)^*\) despite its infinite index set.

The standing hypothesis of Proposition~\ref{prop:WV} and
Corollary~\ref{cor:WV-arm} is \(V\in C^\infty(M)\).  The genericity claim
has the following corrected \(C^k\) form.  For \(k\ge1\) and \(V\in C^k(M)\),
let
\[
L_V:\mathcal H^1(M)\longrightarrow C^{k-1}(M),
\qquad L_V(h)=h(\nabla V),
\qquad W(V):=\Phi(\ker L_V).
\]
Then
\[
\{V\in C^k(M):W(V)=\{0\}\}
\]
is open and dense in \(C^k(M)\).  Here is the short proof, included to
avoid relying on the predecessor wording that placed a \(C^\infty\)
subset inside the \(C^k\) topology.  If \(b_1(M)=0\) the claim is immediate.
Otherwise fix a harmonic basis \(h_1,\ldots,h_b\).  Injectivity of \(L_V\)
is equivalent to the existence of points \(p_1,\ldots,p_b\) for which
\[
\det\!\bigl[h_i(\nabla V)(p_j)\bigr]_{i,j=1}^b\neq0.
\]
This determinant depends continuously on the first jet of \(V\), proving
openness.  For density, the evaluation covectors
\(h\mapsto h_p(v)\), with \(p\in M\) and \(v\in T_pM\), span
\(\mathcal H^1(M)^*\): a harmonic form annihilating all of them is zero.
Choose distinct \(p_j\) and \(v_j\in T_{p_j}M\) whose evaluation covectors
form a basis, and smooth bump functions \(\psi_j\) with disjoint supports
and \(\nabla\psi_j(p_j)=v_j\).  For
\(V_s=V_0+\sum_js_j\psi_j\), the displayed determinant is a polynomial in
\(s=(s_1,\ldots,s_b)\) with non-zero top-degree coefficient.  Arbitrarily
small generic \(s\) therefore makes it non-zero, proving density in every
\(C^k\) topology with \(k\ge1\).

This is a realizability obstruction on a specified compact admissible
class, not a collapse of the two structural types.  Note its direction:
for generic \(V\) the compact admissible topological-only sector is
\emph{empty}, so
positive instances such as Example~\ref{ex:WV-torus} must be exhibited by
construction and never expected generically.

\section{Identity, cortical waves, and the dissociation test}

The topological channel carries identity only on the declared
identity-transfer domain of \cite{structural}.  It does not follow that
every non-zero winding in nature carries that semantic role.  The
correct empirical move is to transport the invariant type first and test
the proposed role separately.

For each condition \(c\in\{\mathrm{pre},\mathrm{post}\}\), let
\(U_c\) be the registered cortical domain and let
\[
z_B^{(c)}:=\mathcal H_B(x^{(c)}),\qquad
a^{(c)}:=|z_B^{(c)}|,\qquad
\phi_c:=z_B^{(c)}/|z_B^{(c)}|
\]
be the registered band-limited analytic signal, its amplitude and its
phase on the phase domain \(D_c:=\{a^{(c)}>a_{\min}\}\), where the
pointwise floor \(a_{\min}>0\) is fixed before inspection.  Write
\(d\vartheta\) for the angular \(1\)-form normalized by
\(\oint_{S^1}d\vartheta=2\pi\).  For a piecewise \(C^1\) loop
\(\gamma_c\subset D_c\),
\[
n_{\gamma_c}
  :=\frac{1}{2\pi}\oint_{\gamma_c}\phi_c^*(d\vartheta)\in\mathbb Z.
\]
This is the same invariant type---degree/winding around a puncture---as
the topological readout above.  It is not literally the same physical
carrier or the same identity certificate.

For an executable discrete certificate, sample a complete oriented loop at
\(p_0,\ldots,p_{N-1}\), obtaining noisy observed values
\(z_i=z_B(p_i)\).  Independently register one cyclically glued,
piecewise-\(C^1\), non-vanishing reference reconstruction \(\zeta\) on the
oriented edges \(e_i:p_i\to p_{i+1}\).  The observed values need not be exact
endpoint values of \(\zeta\).  Let \(\widetilde\vartheta_i^*\) be the chosen
continuous phase lift of \(\zeta|_{e_i}\), write
\[
\delta_i^*:=\widetilde\vartheta_i^*(1)-\widetilde\vartheta_i^*(0),
\qquad \theta_i:=\arg z_i,
\]
and set cyclically
\[
\Delta_i:=\operatorname{Arg}
 \exp\!\bigl(\mathrm i(\theta_{i+1}-\theta_i)\bigr)\in(-\pi,\pi],
\qquad
\widehat n:=\frac1{2\pi}\sum_{i=0}^{N-1}\Delta_i.
\]
Register a within-condition reference-reconstruction/edge-lift certificate, an
independently justified mesh/regularity bound \(b_\phi\), a certified per-edge
\emph{linear, branch-coupled} observation error
\(\varepsilon_{\mathrm{step}}\ge0\), an ambiguity margin
\(0<\kappa<\pi\), and an integer tolerance
\(0\le\varepsilon_{\mathbb Z}<1/2\).  The certificate records every real lift
increment \(\delta_i^*\), pairs it with the observed principal increment on
the declared branch, and either defines \(\zeta\) as the operational phase
target or supplies a within-condition degree-preserving bridge from \(\zeta\)
to \(\phi|_\gamma\).  Thus the next absolute value is on \(\mathbb R\), not a
circular distance:
\[
|\delta_i^*|\le b_\phi,\qquad
|\Delta_i-\delta_i^*|\le\varepsilon_{\mathrm{step}}
\quad\text{for every }i.
\]
The numerical amplitude ledger records
\[
a_{\mathrm{obs,min}}:=\min_i|z_i|,
\qquad
a_{\mathrm{ref,min}}:=\inf_{t\in S^1}|\zeta(t)|.
\]
A winding value is certified only if both recorded minima strictly exceed the
predeclared positive floor \(a_{\min}\), the certificates cover every edge, and
\[
b_\phi+\varepsilon_{\mathrm{step}}<\pi-\kappa,\qquad
N\varepsilon_{\mathrm{step}}<\pi,\qquad
\max_i|\Delta_i|\le b_\phi+\varepsilon_{\mathrm{step}},\qquad
\bigl|\widehat n-\operatorname{round}(\widehat n)\bigr|
 \le\varepsilon_{\mathbb Z}.
\]
The reference reconstruction satisfies
\(n_{\mathrm{ref}}=(2\pi)^{-1}\sum_i\delta_i^*\in\mathbb Z\), while
\[
|\widehat n-n_{\mathrm{ref}}|
 \le \frac{N\varepsilon_{\mathrm{step}}}{2\pi}<\frac12.
\]
Consequently
\(n_{\mathrm{ref}}=\operatorname{round}(\widehat n)\); when the declared
degree-preserving bridge is present this is also \(n_\gamma\).  The numerical
reference steps, their branch coupling to the observations, the true-step
bound and the aggregate error gate are all load-bearing: an integer residual
by itself cannot exclude a reference/observation full-turn mismatch.  Failure
of any certificate gives non-applicability, not zero winding.

Cross-condition comparison uses an orientation-preserving registered
correspondence \(T_j:\gamma_j^{\mathrm{pre}}\to
\gamma_j^{\mathrm{post}}\), together with endpoint-specific annular
phase-domain and singularity-tracking certificates.  It does \emph{not}
assume a homotopy from \(\phi_{\mathrm{pre}}\) to
\(\phi_{\mathrm{post}}\) through non-vanishing phase maps: such a homotopy
would force equality of the two degrees.  If a time-resolved physical
interpolation is studied, a degree change instead requires a zero-amplitude
or singularity crossing in the intervening world tube; that crossing is a
candidate mechanism to locate, not a premise to exclude.

The Science study of brain-wide rotating waves \cite{ye} analyzes
rotating phase waves primarily in the \(2\text{--}8\) Hz band, reports
brain-wide coordination and behavioural recruitment, and finds that
disrupting circular somatosensory wiring reduces the waves.  It treats
rotation and oscillation amplitude together; it does not introduce the
cohomological identity semantics or the typed separation above.

For an empirical test, predeclare a spatial locus \(\mathcal U\), time
window, calibration, band \(B=[2,8]\,\mathrm{Hz}\), and scalar amplitude
functional \(\mathcal A_{\mathcal U,B}\) applied to the same registered
analytic-signal magnitude:
\[
A_{\mathcal U,B}^{(c)}
  :=\mathcal A_{\mathcal U,B}\!\left(a^{(c)}\right)
  \in\mathbb R_{\ge0}.
\]
A different amplitude source requires an explicit observation bridge and
a validation that the two estimators do not force winding and amplitude to
co-vary.  Predeclare a calibrated detection floor \(A_{\min}>0\) and an
observed-record equivalence margin
\[
0\le\delta_A<A_{\min}.
\]
The instrument calibration must resolve differences at that margin; if its
uncertainty is larger, the arm is non-applicable rather than licensed to
widen \(\delta_A\).  This criterion concerns the registered observed
records.  It is not by itself a confidence statement about equality of
latent population amplitudes.

\begin{hypothesis}[Empirical dissociation criterion]\label{hyp:empirical}
Let \(\{(\gamma_j^{\mathrm{pre}},\gamma_j^{\mathrm{post}},T_j)\}\) be a
preregistered endpoint loop family, and let the evaluated loop be named in
advance or selected under a preregistered multiplicity plan.  Require the complete observed/reference reconstruction, branch-coupled
edge-lift and degree-preserving discrete certificate above in both
conditions, using the same registered analytic-signal, loop, calibration
and sampling pipeline, and
require endpoint annular phase-domain and singularity-tracking certificates.
Under those premises, a loop-local empirical realization of the separation
has, for the evaluated loop,
\[
\begin{gathered}
|n_{\gamma_j}^{\mathrm{pre}}|>0,\qquad
n_{\gamma_j}^{\mathrm{post}}=0,\\
\min\!\left\{A_{\mathcal U,B}^{(\mathrm{pre})},
             A_{\mathcal U,B}^{(\mathrm{post})}\right\}\ge A_{\min},\\
\left|A_{\mathcal U,B}^{(\mathrm{post})}
      -A_{\mathcal U,B}^{(\mathrm{pre})}\right|
   \le\delta_A.
\end{gathered}
\]
This is observed-record \(\mathcal U\)-amplitude-preserved loss of winding
on that loop.  It is loop-local in winding and \(\mathcal U\)-local in
amplitude; only when \(\mathcal U\) is a registered neighbourhood of
\(\gamma_j\) is it a loop-local amplitude/winding dissociation.  A
singularity lying on an endpoint loop makes its winding undefined; recording
that case as zero would manufacture the dissociation from an artefact.

Collapse on one loop bounds nothing on the remaining periods.  If the
classes \(\{[\gamma_j]\}\) generate a registered monitored subspace
\(H_{\mathrm{mon}}\subseteq H_1(U_{\mathrm{post}};\mathbb R)\) and every
post-condition generator has zero winding, then the post-condition
cohomology class vanishes only on \(H_{\mathrm{mon}}\).  Full triviality in
\(H^1_{\dR}(U_{\mathrm{post}})\) follows only when
\(H_{\mathrm{mon}}=H_1(U_{\mathrm{post}};\mathbb R)\); at least one
pre-condition generator must have non-zero winding for a sector collapse.
The stronger \emph{survives-but-dead} description additionally requires
(i) an independently declared viability or functional-continuation
variable to be preserved and (ii) a predeclared bridge validating the
monitored winding as an identity-transfer witness in the target system.
Neither premise is supplied by the amplitude/winding dissociation alone.
\end{hypothesis}

For the dual local-only arm, register one loop \(\gamma\), a neighbouring
region \(\mathcal U_\gamma\), a spatiotemporal window, an oriented
orthonormal frame and units, the derivative/smoothing/mask pipeline, a
same-carrier observation bridge, and a multiplicity plan.  At \(N\ge1\)
uniformly weighted registered samples, let \(J_i\) be the Jacobian of the
same local vector-field estimator and define the coherent signed-spin score
\[
r_i:=\frac12\bigl((J_i)_{21}-(J_i)_{12}\bigr),
\qquad
R_{\mathrm{loc}}:=\left|\frac1N\sum_{i=1}^N r_i\right|.
\]
This estimator does not explicitly consume the registered oscillation-
amplitude functional.  Its mask, derivative scale and reliability may still
depend on signal quality, so empirical amplitude invariance would require a
separate validation.  Oppositely oriented subregions can cancel, so it tests
this registered mean
readout rather than the existence of \(d\omega_S\ne0\) somewhere.  Since
\(R_{\mathrm{loc}}=|\bar r|\) is folded, the procedure is a two-sided test
of the signed mean and is silent on handedness.  Run the
identical full pipeline---including sample count, window, mask, frame,
smoothing, and nuisance fitting---on \(m\) preregistered no-rotation
calibration units, obtaining \(R_1^{(0)},\ldots,R_m^{(0)}\).  Because the
reference distribution is built at the level of whole units, spatial
autocorrelation among the \(N\) within-unit samples does not invalidate the
procedure; the effective null sample size is \(m\), not \(N\).  Before the
test, fix
\[
0<\alpha_{\mathrm{loc}}<1,
\qquad R_{\min}>0,
\qquad \frac1{m+1}\le\alpha_{\mathrm{loc}},
\]
where \(R_{\min}\) is justified by instrument resolution or a declared
minimum practical effect.  The last condition is necessary and sufficient
for the rank grid to contain a value at most
\(\alpha_{\mathrm{loc}}\), since its minimum is \(1/(m+1)\); it does not
guarantee rejection in realized data, where ties or a non-extreme observed
score can keep \(p_{\mathrm{loc}}\) larger.  Now assume the registered \emph{design} premise---logically
independent of the sharp no-rotation null, and not testable from these
data---that the experimental unit and the \(m\) calibration units are
exchangeable draws from a common data-generating distribution, matched in
signal-to-noise, coverage, mask geometry and nuisance structure, so that
under the sharp null the vector
\((R_{\mathrm{loc}},R_1^{(0)},\ldots,R_m^{(0)})\) is permutation-invariant.
Running the identical pipeline controls pipeline heterogeneity, not unit
heterogeneity; unlike a within-unit randomization null, this between-unit
reference distribution is assumed rather than constructed, so the
validity of the arm is exactly as good as that matching.  Under that
premise, use the conservative rank value
\[
p_{\mathrm{loc}}
:=\frac{1+\#\{k:R_k^{(0)}\ge R_{\mathrm{loc}}\}}{m+1}.
\]

\begin{hypothesis}[Calibrated local-only criterion]
\label{hyp:local-only}
A loop/region-local empirical realization of the local-only sector is
declared exactly when a winding estimate satisfies the phase-domain,
registered-loop, observed/reference reconstruction, branch-coupled
edge-lift, aggregate-error, and
sub-\(\pi\) no-alias gates and obeys
\[
n_\gamma=0,
\qquad R_{\mathrm{loc}}\ge R_{\min},
\qquad p_{\mathrm{loc}}\le\alpha_{\mathrm{loc}}.
\]
The use of \(\ge\) in the null exceedance count makes the rank test
conservative under ties.  Here \(R_{\min}\) is a practical decision floor;
the displayed rule alone is not a confidence-bound claim that the latent
effect exceeds \(R_{\min}\).  Invalid winding, same-carrier bridge,
provenance, calibration, cross-unit matching, joint-exchangeability, or
multiplicity premises
are non-applicability, not negative evidence.  When the premises are valid
but any displayed gate fails, the verdict is \emph{no local-only
detection} for the registered loop/region/window.  A family-wide statement
additionally requires generator completeness and its preregistered
multiplicity correction.  Neither outcome refutes
Proposition~\ref{prop:independence}.
\end{hypothesis}

\paragraph{Falsifiable population formulation.}
The two hypotheses above are unit-level declaration rules.  To attach a
falsifiable population claim to either arm
\(q\in\{\mathrm{wind},\mathrm{loc}\}\), preregister a source-population law
\(\nu_q\), an outcome-independent eligibility event \(\mathsf E_q\) with
\(\Pr_{\nu_q}(\mathsf E_q=1)>0\), and the eligible-unit law
\[
\mu_q:=\nu_q(\,\cdot\mid\mathsf E_q=1).
\]
Also predeclare a minimum true-realization prevalence
\(\pi_q\in(0,1)\), a detection-sensitivity lower bound
\(s_q\in(0,1]\), and a test level \(\alpha_q\in(0,1)\).  With
\(\mathsf R_q\) denoting a true realization and \(\mathsf D_q\) its
registered detection, the conjecture is explicitly conditional on
eligibility:
\[
\Pr_{\nu_q}(\mathsf R_q\mid\mathsf E_q=1)
 =\Pr_{\mu_q}(\mathsf R_q)\ge\pi_q,
\qquad
\Pr_{\nu_q}(\mathsf D_q\mid\mathsf R_q,\mathsf E_q=1)
 =\Pr_{\mu_q}(\mathsf D_q\mid\mathsf R_q)\ge s_q.
\]
For \(N_q\) independent draws from \(\mu_q\) and \(K_q\) detections, it is
disconfirmed at the registered eligible-unit level when
\[
\sum_{r=0}^{K_q}\binom{N_q}{r}
 (\pi_qs_q)^r(1-\pi_qs_q)^{N_q-r}\le\alpha_q.
\]
In particular, zero detections require
\((1-\pi_qs_q)^{N_q}\le\alpha_q\).  Clustered or dependent units require a
preregistered dependence-valid replacement bound.  Non-applicable units
are reported but are not negative observations, and post-outcome exclusion
or sampling from a law other than the declared \(\mu_q\) invalidates the
arm.  This probabilistic claim is distinct from the formal existence of the
mathematical channel sectors.

The formal theory predicts that this sector is not type-forbidden.
Whether a particular cortex, task, and estimator class realizes it remains
an empirical question.

\section{Falsification surface and limits}

\begin{enumerate}[leftmargin=2.2em]
  \item Proposition~\ref{prop:independence} is established by two explicit
        fields on one flat compact torus under the canonical Hodge projection,
        so the route to falsification is an arithmetic, Hodge-projection or
        hypothesis error in those constructions.  The zero--zero exact form
        noted afterwards is not a non-gradient carrier witness.
  \item Proposition~\ref{prop:WV} is executable through
        Corollary~\ref{cor:WV-arm}: after the compactness, harmonic-basis,
        class-coordinate, evaluation-span, and exact-arithmetic
        (symbolic or rational) certificates are supplied,
        the observable is \(\rho_V(c)=\lVert Ea\rVert_\infty\) and the
        threshold is exactly zero.  A missing span certificate makes a
        sampled zero non-applicable as a membership certificate; a
        floating-point pipeline places the instance outside this arm
        rather than merely degrading it.
  \item Proposition~\ref{prop:norm} concerns the same one-form on the
        same loop.  Applying it across unmatched instruments is
        non-applicability, not a theorem failure.
  \item The empirical bridge is weakened if winding cannot be estimated
        robustly under registered-loop transport and singularity tracking.
        It is \emph{non-applicable}, not merely weakened, if the
        phase-domain, registered-loop, observed/reference reconstruction and
        branch-coupled edge-lift coupling,
        aggregate-error or sub-\(\pi\) no-alias gates fail,
        or if amplitude and winding are algorithmically forced to co-vary
        by the estimator: the latter is a same-carrier bridge failure.
  \item Hypotheses~\ref{hyp:empirical} and
        \ref{hyp:local-only} are unit-level declaration rules.  The
        falsifiable object is the registered eligible-unit prevalence claim
        \(\Pr_{\mu_q}(\mathsf R_q)\ge\pi_q\), together with its declared
        sensitivity, eligibility-conditioning and sampling premises.  Its
        falsifier measures the registered detection count \(K_q\) and uses
        the displayed lower-tail threshold; it does not treat finite
        non-observation as the logical negation of existence.  Invalid phase,
        provenance, unit matching, eligibility law or exchangeability
        premises are non-applicability.
  \item Failure to observe either dissociation disconfirms the registered
        population claim only when its prevalence/sensitivity/power bound
        crosses the preregistered level.  Otherwise it is no detection, and
        in no case does it refute Proposition~\ref{prop:independence}.
  \item The genericity statement proved in
        Section~\ref{sec:compact} is itself a limit on the programme: the
        compact admissible topological-only sector is empty for generic \(V\),
        so
        Hypothesis~\ref{hyp:empirical} hunts a non-generic sector on that
        branch.
  \item A preserved amplitude without an independent continuation
        variable is not evidence of survival, and winding loss without a
        validated identity-witness bridge is not evidence of identity
        death.  A zero winding on a probe domain with trivial \(H^1\) is
        not evidence that every form of identity is absent.
\end{enumerate}

\section{Conclusion}

Structural spin is one non-gradient layer with two formal channels:
\[
\Sigspin(F;P_{\mathrm{cl}})=([\eta_S],d\omega_S),
\]
with the selector suppressed only on the compact Hodge branch.  They are
logically independent over the class of all declared signature data, and
jointly constrained---generically to the point of collapsing the
topological-only sector---under the compact admissibility package.  Period,
path norm/amplitude, and vorticity flux are three ways to read a declared
one-form drawn from that layer, once probes are chosen.  The dynamic
residual is a further ledger variable.  The norm inequality connects a
period to a compatible path norm and nothing more.  The compact
topological-only premise gate is executable, conditional on the registered
basis, evaluation-span and exact-arithmetic certificates, through the exact
residual \(Ea\).  The dual empirical arm returns a three-valued verdict at
its declared resolution---non-applicability, no detection, detection---with
certified zero winding requiring a positive phase-domain floor, an
independent sub-\(\pi\) true-step bound, a numerical reference edge-lift
ledger with branch-coupled observation errors, an aggregate error below half
a turn, an integer tolerance, a positive observed/reference amplitude floor,
and complete loop registration.  Once these types are kept separate, the central
experimental bet becomes both stronger and cleaner: the amplitude statistic
from the same registered analytic signal can remain above its detection floor
and within its observed-record equivalence margin while winding is lost on
the registered loop---or, under the generator conditions of
Hypothesis~\ref{hyp:empirical}, on the monitored homology subspace.  Full
cohomology-class collapse requires generators of all \(H_1\); population
disconfirmation requires the separately registered eligible-unit
prevalence/sensitivity bound.  Identity death still requires the independent
transfer and viability premises that give the winding witness its meaning.

\begin{thebibliography}{9}

\bibitem{structural}
M.~Barziankou,
\emph{Structural Spin as a Forgetful Separation over Dissipative
Dynamics, with a Cohomological Obstruction to Identity Transfer},
Navigational Cybernetics~2.5 (2026).
\href{https://doi.org/10.17605/OSF.IO/94GWQ}
{doi:10.17605/OSF.IO/94GWQ}.

\bibitem{operational}
M.~Barziankou,
\emph{Operational Spin: A Physical Projection of UTAM-Directed Will into
Adaptive Media} (2025).
\href{https://doi.org/10.5281/zenodo.17938800}
{doi:10.5281/zenodo.17938800}.

\bibitem{bridge}
M.~Barziankou,
\emph{Spin-Channel Bridge and Core-Reduction: A Categorical Foundation}
(2026).
\href{https://doi.org/10.17605/OSF.IO/EAUD5}
{doi:10.17605/OSF.IO/EAUD5}.

\bibitem{nc25}
M.~Barziankou,
\emph{Navigational Cybernetics~2.5}.
\href{https://doi.org/10.17605/OSF.IO/NHTC5}
{doi:10.17605/OSF.IO/NHTC5}.

\bibitem{ye}
Z.~Ye et al.,
``Brain-wide topographic coordination of rotating waves'',
\emph{Science} \textbf{392} (2026), eadx1369.
\href{https://doi.org/10.1126/science.adx1369}
{doi:10.1126/science.adx1369}.

\bibitem{bott-tu}
R.~Bott and L.~W.~Tu,
\emph{Differential Forms in Algebraic Topology},
Springer, 1982.

\bibitem{warner}
F.~W.~Warner,
\emph{Foundations of Differentiable Manifolds and Lie Groups},
Springer, 1983, Ch.~6.

\end{thebibliography}

\end{document}
